% !TeX root = ../../arithmetic-and-computation.tex
\section{Individual self-dual models and lost automorphisms}
These constructions establish the local obstructions and special-fiber
alternatives. Paper I uses one common nonfree subgroup for its global
family; an arbitrary individual choice is not silently substituted for it.

\begin{proposition}[Principally polarized model at the $C$-vertex]\label{prop:Cppisogenous}
The normalized $C$-vertex admits an $\OL$- and $\mathcal O_K$-stable
self-dual overlattice of index
\[
 \operatorname{disc}(L)=453789
\]
for $K=\Ql(\sqrt{-3})$.  Hence it is $\OL$-linearly isogenous of degree
$453789$ to a principally polarized sixfold $A_C^{\mathrm{pp}}$ with the same
rational CM Hodge structure, signature, discriminant, and Hodge-conjecture
conclusion.
\end{proposition}
\begin{proof}
Write $\Sigma_C=(I+2C_t^{14})^{\mathsf T}$.  On the discriminant module
\[
 \Lambda^\#/\Lambda\cong
 (\OL/\mathfrak p_3)^2\oplus(\OL/\mathfrak p_7^5)^2
\]
one has $\Sigma_C\equiv0\pmod{\mathfrak p_3}$.  Choose any residue line on that
part.  On the $\mathfrak p_7^5$-part, $\Sigma_C^2=4I$ and
$\operatorname{tr}(\Sigma_C)=0$, so the idempotents
\[
 \frac{\Sigma_C+2I}{4},\qquad \frac{2I-\Sigma_C}{4}
\]
give rank-one eigensummands.  Either is isotropic because
$\Sigma_C^\dagger=-\Sigma_C$.  The CRT product has order $3^3 7^5$ and is therefore
maximal isotropic.  The same argument as for $A$ and $B$ proves
$\mathcal O_K$-stability, self-duality, preservation of all rational Hodge
data.
\end{proof}

\begin{corollary}[The principal $C$-model loses the full vertex automorphism]\label{cor:noCstablepp}
No principally polarized self-dual overlattice of the standard lattice can
retain the full integral automorphism acting as $C_t^{\mathsf T}$ on homology.
\end{corollary}
\begin{proof}
Stability under $C_t^{\mathsf T}$ and its inverse implies stability under
$tI$, so the order-$27$ subgroup is an $\mathbb F_{27}$-line. Modulo
$\mathfrak p_3$, the characteristic discriminant of $C_t^{\mathsf T}$ is
$t^2-1$, and exact arithmetic in
$\mathbb F_{27}=\mathbb F_3[t]/(t^3-t^2+t+1)$ gives
\[
 (t^2-1)^{13}=-1.
\]
Thus the characteristic polynomial is irreducible over $\mathbb F_{27}$ and
$C_t$ has no invariant residue-field line.
\end{proof}

\begin{proposition}[Principally polarized model at the $A$-vertex]\label{prop:Appisogenous}
Let $\Lambda=\mathcal O_L^2$ be the homology lattice at the normalized $A$-vertex and
$\Lambda^\#=\mathfrak D_L^{-1}\Lambda$ its dual for $E_1$.  There is an
$\mathcal O_L$- and $\mathcal O_K$-stable maximal isotropic subgroup
\[
 H_A\subset\Lambda^\#/\Lambda,\qquad K=\mathbb Q(\sqrt{-3}),
\]
whose inverse image $\Lambda_A^{\mathrm{pp}}$ is self-dual for $E_1$.  Hence
\[
 [\Lambda_A^{\mathrm{pp}}:\Lambda]=\operatorname{disc}(L)=453789
\]
and the normalized $A$-vertex sixfold is $\mathcal O_L$-linearly isogenous of degree
$453789$ to a principally polarized abelian sixfold $A_A^{\mathrm{pp}}$.
The rational CM Hodge structure, the $K$-signature $(3,3)$, the Hermitian
discriminant class $-1$, and the divisor-generation of the Hodge algebra are unchanged.
In particular the Hodge conjecture holds for $A_A^{\mathrm{pp}}$ in every codimension.
\end{proposition}
\begin{proof}
As in Proposition~\externalref{packet}{prop:ppisogenous},
\[
 \Lambda^\#/\Lambda\cong
 (\mathcal O_L/\mathfrak p_3)^2\oplus
 (\mathcal O_L/\mathfrak p_7^5)^2,
 \qquad
 \mathfrak D_L=\mathfrak p_3\mathfrak p_7^5.
\]
Let $\Sigma_A$ be the integral operator representing $\sqrt{-3}$ in the proof of
Theorem~\externalref{packet}{thm:Ahodge}.  Exact reduction gives
\[
 \Sigma_A\equiv0\pmod{\mathfrak p_3}.
\]
Thus any one-dimensional $\mathbb F_{27}$-subspace of the $\mathfrak p_3$-part is
$\Sigma_A$-stable; fix one and call it $H_{A,3}$.

On the $\mathfrak p_7^5$-part one has $\Sigma_A^2=-3I=4I$ and
$\operatorname{tr}(\Sigma_A)=0$.  Since $4$ is a unit, the idempotents
\[
 e_{A,+}=\frac{\Sigma_A+2I}{4},\qquad
 e_{A,-}=\frac{2I-\Sigma_A}{4}
\]
split the rank-two local module into two rank-one direct summands.  Put
$H_{A,7}=\operatorname{im}(e_{A,+})$.  The relation
$\Sigma_A^\dagger=-\Sigma_A$ gives, for $v,w\in H_{A,7}$,
\[
 0=E_1(\Sigma_Av,w)+E_1(v,\Sigma_Aw)=4E_1(v,w),
\]
so $H_{A,7}$ is isotropic.  Therefore
\[
 H_A=H_{A,3}\oplus H_{A,7}
\]
is isotropic and has order
\[
 |H_A|=3^3 7^5=\sqrt{|\Lambda^\#/\Lambda|},
\]
hence is maximal isotropic.  Its inverse image is self-dual.

Exact power-basis reduction verifies $U_A=(\Sigma_A-I)/2\in M_2(\OL)$; one has $U_A^2+U_A+I=0$ and
$\mathcal O_K=\mathbb Z[U_A]$.  On the $\mathfrak p_3$-part,
$U_A\equiv1$, while on the $\mathfrak p_7^5$-part it is a polynomial in
$\Sigma_A$; hence $H_A$ is $\mathcal O_K$-stable.  The rational vector space,
complex structure, $K$-action, and alternating form $E_1$ are unchanged by the
lattice enlargement, so the stated Hodge-theoretic invariants are unchanged.
\end{proof}

\begin{corollary}[The principal $A$-model loses the full vertex automorphism]\label{cor:noAstablepp}
No principally polarized self-dual overlattice
\[
 \Lambda\subset\Lambda'\subset\Lambda^\#
\]
can be stable under the full normalized projective order-$7$ automorphism acting as $A_t^{\mathsf T}$ on homology.
\end{corollary}
\begin{proof}
On the $\mathfrak p_3$-part the subgroup has order $27$. Stability under
$A_t^{\mathsf T}$ and its inverse implies stability under their sum
$(t^3-3t)I=t^3I$ in characteristic $3$. Since $(t^3)^9=t$ in
$\mathbb F_{27}$, the subgroup is an $\mathbb F_{27}$-line.  The characteristic polynomial of $A_t$ is
\[
 X^2-(t^3-3t)X+1.
\]
Its discriminant reduces to $1-t$ modulo
$\mathfrak p_3=(3,t^3-t^2+t+1)$, and exact finite-field arithmetic gives
\[
 (1-t)^{13}=-1\in\mathbb F_{27}.
\]
Thus the discriminant is a nonsquare, the characteristic polynomial is irreducible, and
$A_t$ has no invariant residue-field line.
\end{proof}

